% Contribution fragment for the team's AI Transparency Log.
% Requires enumitem and xurl (or url/hyperref) in the parent document.
% Insert inside the team AI appendix, not as a second standalone report.
\subsection*{Algorithm and Interface-PINN Contribution --- Tianrui Ji}
\label{app:ai_tianrui_ji}

\paragraph{Tools and scope.}
ChatGPT/Codex assisted with code, experiment scripts, LaTeX text, figures and
presentation revisions. I organised and integrated these materials and
requested changes. The assistance extended beyond prose polishing;
independent validation is attributed to the testing member's record.

\paragraph{Audit TJ1: Discretisation, solver accuracy and interface checks.}
\begin{itemize}[leftmargin=1.5em,itemsep=2pt]
    \item \textbf{Task/query:} Develop the algorithm section for the
    variable-diffusion problem, including conservative FDM, two-dimensional
    P1 FEM, sparse solvers and reproducible comparisons.
    \item \textbf{AI-assisted output:} Assembly and solver code, numerical
    tables and figures, and draft interpretations of convergence and
    interface diagnostics.
    \item \textbf{Risk and possible falsifier:} Incorrect boundary signs or
    face coefficients can survive symmetry and residual checks. Substituting
    exact nodal data into selected rows and comparing numerical and analytic
    fluxes can expose these errors.
    \item \textbf{Verification and evidence:} The assembly is explicit in
    \path{project2/algorithm_section/src/discretisation.py}, particularly
    \texttt{assemble\_fdm} and \texttt{assemble\_fem}; the solver interface
    is in \path{project2/algorithm_section/src/solvers.py}. The testing
    member's independent evidence is recorded in
    \path{project2/test/VALIDATION_RUN_2026-11-06.md} and
    \path{project2/test/results/}. It checks manufactured solutions, linear
    solves, matrix spectra and interface diagnostics.
    \item \textbf{Decision and limitation:} Algebraic, discretisation and
    physical flux errors are reported separately. At $c=10$, arithmetic
    averaging reduces nodal error under refinement but retains maximum
    interface-flux error near $0.662$. The report uses the independently
    corrected condition numbers. Conclusions remain benchmark-specific.
\end{itemize}

\paragraph{Audit TJ2: Interface-PINN comparison and training provenance.}
\begin{itemize}[leftmargin=1.5em,itemsep=2pt]
    \item \textbf{Task/query:} Present the interface network as an
    alternative solver for the same layered problem, compare it with FDM
    and FEM, and disclose the architecture and training cost.
    \item \textbf{AI-assisted output:} Two-network code, experiment records
    and initial slides comparing mainly with a single-network PINN. Later
    drafts added three-method evaluation and training plots.
    \item \textbf{Risk and possible falsifier:} Norms and sampling must match.
    Independent tanh networks differ from published shared-parameter
    I-PINNs, and low loss does not establish physical accuracy. Check device
    claims against the environment record and errors against exact data.
    \item \textbf{Verification and evidence:} Configuration and device
    records are in
    \path{project2/I-PINNs/runs/server_interface/config.json} and
    \path{project2/I-PINNs/runs/server_interface/environment.json}.
    Each run directory contains \texttt{training.csv},
    \texttt{metrics.json} and \texttt{checkpoint.pt}; for example,
    \path{project2/I-PINNs/runs/server_interface/interface_c10_seed11/}.
    The independent checkpoint and metric checks are documented in the
    validation record cited above.
    \item \textbf{Decision and limitation:} I requested three-method
    comparison and clarification of joint training. Records identify an
    A100-PCIE-40GB; mean optimisation time at $c=10$ is $231.3$ seconds.
    Only periodic Adam losses and a final L-BFGS loss are saved, not a full
    L-BFGS or error-versus-time trajectory. The revised comparison and plots
    are integrated in Section~\ref{sec:interface_pinn}. The three checkpoints
    used there and their common-sample replay are under
    \path{project2/report/support/interface_pinn/}. Sampled relative
    errors are distinguished from continuous $L^2$ norms; no retraining
    was performed for the integration. A later integration replays three
    single-network checkpoints as a secondary ablation; its doubled
    parameter count in the interface model is disclosed.
\end{itemize}

\paragraph{Responsibility and trusted dependencies.}
Team-owned code specifies PDE data, assembly, boundary treatment and
experiment protocols. SciPy supplies sparse linear algebra; PyTorch supplies
automatic differentiation and optimisers. My contribution and learning
limits are stated in \path{TianruiJiICS.tex}. This fragment does not sign a
declaration for other members.
